\section{Literature Review}

\subsection{LLMs as synthetic respondents}

LLMs offer inexpensive ways to pilot surveys and experimental vignettes. Early work explored demographic conditioning and \emph{algorithmic fidelity} \cite{argyle2023}, classical experimental replication \cite{aher2023}, and simulated economic agents \cite{horton2023}. Their usefulness depends not only on plausible answers but on reproducing human treatment contrasts. Replication studies find some directional agreement alongside exaggerated magnitudes, compressed variation, and weaker correspondence in socially sensitive domains \cite{cui2025,bisbee2024,santurkar2023,rottger2024}. We therefore benchmark LLM treatment contrasts against published human studies rather than assuming that synthetic samples are representative.

\subsection{Social desirability and instruction sensitivity}

Social desirability bias (SDB) is response distortion associated with perceived acceptability. Survey theory distinguishes impression management from self-deceptive enhancement and emphasizes question sensitivity, privacy, and administration mode \cite{crowne1960,paulhus1984,tourangeau2007,krumpal2013}. Indirect elicitation and anonymity can reduce some human response pressures \cite{nederhof1985,horiuchi2022}; system prompts are not established equivalents of these survey methods.

LLM behavior introduces additional explanations. Models may favor socially approved answers, conform to perceived user expectations, or respond to safety-related training and instructions \cite{salecha2024,sharma2023}. Survey prompts and response formats can also change their outputs \cite{tjuatja2024}. A shift away from a socially approved answer therefore need not represent reduced human-like SDB. This motivates studying system-prompt stimuli as experimentally consequential inputs and distinguishing behavioral steering from correspondence with human treatment effects.

\subsection{Reasoning and internal representations}

Chain-of-thought prompting can alter generated outputs \cite{wei2022,kojima2022}, but verbalized reasoning is not a validated record of a model's internal causal process. We distinguish final answers, coded argumentative content, and hidden-state measurements. Annotation can identify associations between expressed arguments and responses, conditional on its reliability. Hidden-state probes offer a separate representational diagnostic; in our design, they concern a different local model and cannot be treated as a direct explanation of API-model behavior.

\subsection*{Research contribution}

This study compares original treatment effects and system-prompt effects across four adapted experimental domains, with published human contrasts as benchmarks. For the two original cases it also examines reasoning annotation and exploratory hidden-state evidence. The extension tests whether prompt sensitivity generalizes across models, normative contexts, and response formats without assuming that movement away from socially desirable answers improves human validity.
